%% AASTeX v7.0.1 (released May 29, 2025)
%% Classfile: aastex701.c0988ls
%% Download from: https://journals.aas.org/aastex-package-for-manuscript-preparation/
\documentclass[preprint]{aastex701-1/aastex701}

\usepackage{amsmath}
\usepackage{amssymb}
\usepackage{graphicx}
\usepackage{placeins} % Added to constrain floats

\begin{document}

\title{Prism: A Physics-Constrained Neural Emulator for Stellar Spectra}

\author[0009-0006-8079-688X]{Aneesh Shastri}
\affiliation{Department of Computational Engineering, Cepheid,
Indian Institute of Technology Hyderabad,
Sangareddy, Telangana 502284, India}
\email{co25btech11004@iith.ac.in}

\begin{abstract}
We present Prism, a physics-constrained generative neural network emulator for stellar spectra
trained directly on real, high SNR observations from the APOGEE DR17 survey.
Unlike prior approaches that rely on synthetic spectral grids, we train our model on
149,966 stars with signal-to-noise ratio $>100$ using stellar parameters and chemical
abundances from the ASPCAP pipeline as labels, spanning a wavelength range of 1514--1694\,nm.

The current repository implements Prism as a two-input masked model: in addition to the
27-dimensional label vector, the network receives a 23-dimensional per-label imputation
mask and is trained against a 3-channel target tensor containing flux, inverse variance,
and a per-star pixel mask. The architecture predicts optical depth through a continuum
branch, a CNO molecular branch, a thermodynamic state branch, and element-specific sparse
atomic experts, then converts to flux via the Beer-Lambert law. The active trainer also
retains a single-input structured ablation used as the comparison model ``Ablation Prism,''
which preserves the sparse atomic projection masks but feeds the full label vector through
its internal branches.

Evaluation in the current codebase is explicitly imputation-aware: stars with imputed core,
CNO, or Fe labels are dropped, while trace-element imputations are handled by masking the
corresponding spectral windows. On the retained 8,938-star subset of the held-out test set,
Prism attains a median reduced $\chi^2_\mathrm{red}=1.718$, competitive with the
pre-trained Payne-family baselines loaded by the evaluation notebook. Prism's defining
advantage is structural rather than purely residual-based: an information-leakage test
returns exactly zero forbidden-region response for all 14 atomic species, while every
baseline exhibits non-zero cross-element leakage. We further validate a fast inverse
pipeline built around a heteroscedastic CNN warm start and XLA-compiled MAP optimization;
on a 500-star stratified sample it achieves median wall-time 0.31\,s per star with
RMSEs of 218\,K in $T_\mathrm{eff}$, 0.107\,dex in $\log g$, and 0.039\,dex in
$[\mathrm{M/H}]$, but much smaller MADs of 10.3\,K, 0.038\,dex, and 0.013\,dex,
respectively, indicating that the inverse fit is usually tight and that the inflated RMSE
is driven by a small number of severe outliers. The resulting repository now supports a
more nuanced conclusion: hard spectral locality yields exact disentanglement and
competitive forward fits, with especially strong behaviour in weak-line regimes such as
metal-poor giants.
\end{abstract}

\keywords{stars: abundances --- stars: fundamental parameters --- methods: data analysis ---
techniques: spectroscopic --- surveys}

%------------------------------------------------------------
\section{Introduction}
\label{sec:intro}
%------------------------------------------------------------

The chemical fingerprint encoded in a stellar photosphere reflects the conditions of the
interstellar medium at the time and place of that star's formation, as well as the integrated
nucleosynthetic enrichment from preceding generations of supernovae and asymptotic giant
branch stars. Since photospheric abundances change negligibly over a stellar lifetime,
high-fidelity abundance measurements from large spectroscopic surveys constitute a powerful
fossil record of Galactic chemical evolution \citep{Freeman2002,BlandHawthorn2016}.

This goal has only recently become tractable as survey facilities deliver high-resolution
($R \gtrsim 20{,}000$) spectra for $\sim 10^5$--$10^6$ stars across all Galactic components
\citep{Majewski2017}. APOGEE (Apache Point Observatory Galactic Evolution Experiment)
in particular has assembled one of the largest homogeneous near-infrared spectroscopic
datasets in existence, with Data Release 17 containing labels derived for hundreds of
thousands of giant stars \citep{Abdurrouf2022}.

Automated pipelines such as ASPCAP (APOGEE Stellar Parameter and Chemical Abundance
Pipeline) have made survey-scale stellar characterization possible, yet physics-based model
spectra are known to be incomplete: one-dimensional model atmospheres, the local
thermodynamic equilibrium approximation, and imperfect atomic and molecular line lists
all introduce systematic residuals \citep{GarciaperezASPCAP2016}. These limitations become
especially acute at low signal-to-noise ratios, where classical $\chi^2$ fitting degrades
rapidly.

Data-driven models offer a complementary approach. The Cannon \citep{Ness2015Cannon}
demonstrated that a quadratic generative model trained on high-fidelity reference spectra
can transfer stellar labels to lower signal-to-noise data with competitive precision. The
Cannon 2 \citep{Casey2016Cannon2} extended this to 17 labels simultaneously by introducing
$\ell_1$ regularization that enforces sparsity in the model's spectral coefficients, yielding
physically interpretable spectral derivatives. The Payne \citep{Ting2019Payne} replaced the
polynomial basis with a multi-layer perceptron trained on grids of synthetic spectra,
achieving high accuracy on APOGEE observations. More recently, TransformerPayne
\citep{TransformerPayne2025} employed large transformer architectures over the same
synthetic-grid paradigm, reporting mean absolute errors of $\sim 10^{-3}$ on held-out
synthetic data.

A common thread in nearly all prior emulators is reliance on synthetic spectral grids
for training. While such grids are noise-free by construction, they encode all the
approximations of the underlying physics model and, critically, are evaluated on real data
that differ systematically from their training distribution. Furthermore, most neural emulators
operate as black boxes: the learned mapping from stellar parameters to predicted flux offers
little interpretable insight into which spectral features are being used, leaving open the
question of parameter cross-talk (entanglement).

In this work, we present a generative neural emulator trained entirely on real APOGEE DR17
spectra. Prism is distinguished by three design principles. First, we predict optical
depth rather than flux directly, using the Beer-Lambert law to recover flux from summed
contributions of physically separated branches. Second, element-specific spectral masks
enforce strict disentanglement: each chemical abundance label is permitted to influence only
those pixels lying within 0.2\,nm of its known atomic lines, as compiled in
\citet{Shetrone2015} and \citet{Smith2021DR17}. Third, because our training set is composed
of real observations with realistic noise, continuum normalization artifacts, and sensor gaps,
the emulator inherits the noise properties of APOGEE data without requiring a synthetic
bridge.

Throughout this work we use the term \emph{label} for any element of the input
parameter vector, including stellar parameters ($T_\mathrm{eff}$, $\log g$, etc.) and
chemical abundances ([$X$/Fe] or [$X$/H]).

%------------------------------------------------------------
\section{Data}
\label{sec:data}
%------------------------------------------------------------

\subsection{APOGEE DR17}
\label{subsec:data_apogee}

We use spectra and labels from APOGEE Data Release 17 \citep{Abdurrouf2022}. The APOGEE
spectrograph covers the $H$-band (1514--1694\,nm excluding sensor gaps) at resolution $R \approx 22{,}500$.
The ASPCAP pipeline performs continuum normalization and spectral stacking before
data release \citep{GarciaperezASPCAP2016}. We mask sensor gaps (pixels with zero flux) throughout all processing steps.

Our working sample comprises 149,966 stars selected to have signal-to-noise ratio
SNR $> 100$ in the stacked, pseudo-continuum-normalized spectra. We adopt stellar parameters and
23 chemical abundance labels from ASPCAP (see Table~\ref{tab:labels}). Six core
stellar parameters include effective temperature $T_\mathrm{eff}$, surface gravity $\log g$,
overall metallicity $[M/\mathrm{H}]$, microturbulent velocity $v_\mathrm{micro}$,
macroturbulent velocity $v_\mathrm{macro}$, and projected rotational velocity $v \sin i$.
The 17 elemental abundances span the CNO group, the $\alpha$-capture, odd-Z, and Fe-peak
families.

\begin{deluxetable}{ll}
\tablecaption{Labels used in model training \label{tab:labels}}
\tablewidth{0pt}
\tablehead{\colhead{Group} & \colhead{Labels}}
\startdata
Core parameters & $T_\mathrm{eff}$, $\log g$, $[M/\mathrm{H}]$, $v_\mathrm{micro}$,
                  $v_\mathrm{macro}$, $v\sin i$ \\
CNO group        & $[\mathrm{C/Fe}]$, $[\mathrm{N/Fe}]$, $[\mathrm{O/Fe}]$ \\
Iron peak        & $[\mathrm{Fe/H}]$, $[\mathrm{Mn/Fe}]$, $[\mathrm{Ni/Fe}]$,
                  $[\mathrm{Cr/Fe}]$, $[\mathrm{V/Fe}]$, $[\mathrm{Co/Fe}]$ \\
$\alpha$-capture & $[\mathrm{Mg/Fe}]$, $[\mathrm{Si/Fe}]$, $[\mathrm{Ca/Fe}]$,
                  $[\mathrm{Ti/Fe}]$, $[\mathrm{S/Fe}]$ \\
Odd-Z            & $[\mathrm{Al/Fe}]$, $[\mathrm{K/Fe}]$, $[\mathrm{Na/Fe}]$
\enddata
\end{deluxetable}

\subsection{Preprocessing and Label Imputation}
\label{subsec:preproc}

Missing label values occur in the dataset because ASPCAP flags individual
abundance fits when signal is insufficient or assigns invalid placeholder values. 
We apply robust filtering logic (flagging values $< -100$ or possessing non-zero ASPCAP error flags)
to accurately identify all missing dimensions. 
Rather than discard affected stars, we recover
missing values via iterative imputation using a Bayesian Ridge regressor that learns
inter-label correlations (e.g., the tight correlation between $[\mathrm{Fe/H}]$ and
$[M/\mathrm{H}]$) and predicts missing entries from available ones. To strictly prevent information leakage, the imputer is fitted entirely on the training set and subsequently applied to validation and test data without refitting. After imputation,
we convert all [$X$/Fe] ratios to [$X$/H] by adding the imputed $[\mathrm{Fe/H}]$
value. Details on missing-value fractions and the imputation procedure are available in 
our project repository. A summary of the parameter missing value percentages is presented in Table~\ref{tab:parameter_missing_values}.

\begin{deluxetable}{lrclr}
\tablecaption{Missing Value Percentages \label{tab:parameter_missing_values}}
\tablewidth{0pt}
\tablehead{
    \colhead{Parameter} & \colhead{\%} & \colhead{\hspace{0.5cm}} &
    \colhead{Parameter} & \colhead{\%}
}
\startdata
TEFF   & 0.00\%  && CA\_FE & 1.01\%  \\
LOGG   & 0.00\%  && TI\_FE & 10.14\% \\
M\_H   & 0.00\%  && S\_FE  & 5.17\%  \\
VMICRO & 0.00\%  && AL\_FE & 4.91\%  \\
VMACRO & 0.00\%  && MN\_FE & 8.29\%  \\
VSINI  & 0.00\%  && NI\_FE & 2.05\%  \\
C\_FE  & 0.55\%  && CR\_FE & 5.18\%  \\
N\_FE  & 10.34\% && K\_FE  & 3.24\%  \\
O\_FE  & 0.50\%  && NA\_FE & 4.33\%  \\
FE\_H  & 0.47\%  && V\_FE  & 2.51\%  \\
MG\_FE & 0.47\%  && CO\_FE & 13.81\% \\
SI\_FE & 0.46\%  && \multicolumn{2}{c}{} \\
\enddata
\end{deluxetable}

We supplement the raw labels with four engineered features to improve the model's ability to
decode stellar physics:
\begin{enumerate}
\item $5040/T_\mathrm{eff}$ -- encodes the Boltzmann exponential dependence of line opacity
      on excitation temperature; 
\item $[\mathrm{C/H}] - [\mathrm{O/H}]$ -- the C/O ratio, which governs the balance
      between CO molecule formation and free atomic C and O;
\item $v_\mathrm{broad} = \sqrt{v_\mathrm{macro}^2 + (v\sin i)^2}$ -- a single
      combined broadening term;
\item $\log P_e = 0.5[\mathrm{M/H}]+0.5\log{g}$ -- a first-order approximation of electron pressure derived from
      metallicity and surface gravity, which modulates continuous opacity.
\end{enumerate}

We apply $z$-score normalization to all labels (raw and engineered) using training-set statistics.
Spectral fluxes are used as delivered by ASPCAP (already pseudo-continuum normalized).

The current training pipeline augments each star with two additional mask objects derived
from the imputation flags. First, a 23-dimensional \texttt{bad\_mask} records which input
labels were imputed. Second, a per-star 8575-pixel mask suppresses wavelengths associated
with imputed C, N, O, or trace-element abundances using \texttt{apogee\_mask.npy} and the
APOGEE \texttt{.filt} windows. These masks are serialized into the TFRecord files used by
the active trainer, allowing both the loss and the architecture to treat imputed labels
explicitly rather than relying on naive evaluation against filled-in targets.

%------------------------------------------------------------
\section{Model Architecture}
\label{sec:model}
%------------------------------------------------------------

\subsection{Optical-depth Formulation}
\label{subsec:tau}

Rather than regressing directly to continuum-normalized flux, our model predicts the
effective optical depth $\tau(\lambda)$ and converts to flux via the Beer-Lambert law:
\begin{equation}
F(\lambda) = k \, e^{-\tau(\lambda)},
\label{eq:beer}
\end{equation}

where $k = k(\lambda)$ is a wavelength-dependent scaling factor produced by the continuum branch. We chose a scaling factor instead of another optical depth term to address ASPCAP's continuum normalization artifacts (which occasionally produce normalized fluxes slightly above unity near absorption features).

The total optical depth is additive in contributions from
physically distinct components:
\begin{equation}
\tau(\lambda) =  \tau_\mathrm{CNO}(\lambda)
               + \sum_X \tau_X(\lambda),
\label{eq:tau_sum}
\end{equation}
where the sum runs over all atomic species $X$ in Table~\ref{tab:labels} (excluding the
CNO group), and each term is produced by a dedicated network branch.

\subsection{Network Branches}
\label{subsec:branches}

\subsubsection{Continuum Branch}
The current trainer replaces the earlier upsampling CNN continuum with a rigid
Chebyshev-basis projection. A small multilayer perceptron maps the selected atmospheric
inputs to 20 polynomial coefficients, and a custom \texttt{ChebyshevContinuum} layer
projects those coefficients directly onto the 8575-pixel wavelength grid. This produces a
smooth pseudo-continuum by construction while sharply limiting the branch's ability to fit
narrow atomic structure, which is delegated to the molecular and atomic opacity branches.
Architecturally, this is a meaningful change from the older learned upsampling stack:
the continuum is now constrained to live in a low-order polynomial basis, making it both
more interpretable and less able to absorb narrow line-like residuals that should instead
be explained by the opacity branches.

\subsubsection{CNO (Molecular) Branch}
The molecular branch remains dense and global. In the two-input Prism model it receives a
concatenation of selected atmospheric, rotational, and CNO quantities, and the codebase
supports optional per-sample \texttt{GradientGate} suppression of imputed C, N, and O
inputs. Two dense layers of widths 256 and 1024 predict $\tau_\mathrm{CNO}$ on the full
8575-pixel grid. As in earlier versions, we leave this branch unmasked because blended
molecular opacity is intrinsically non-local in wavelength space. The purpose of
\texttt{GradientGate} is not to alter the forward value of an imputed label, but to block
back-propagated dependence on that label for the affected sample. In effect, the notebook
lets the model consume the filled value for numerical continuity while preventing the CNO
branch from learning spurious gradients from labels that were not actually measured.

\subsubsection{Thermodynamic State Branch}
The thermodynamic state branch distills all global and CNO parameters into a compact
latent state vector. In the current implementation the state branch uses two hidden layers
of widths 32 and 16 followed by a 16-dimensional linear state vector. We pass this state vector to the chemical expert
branch rather than raw label values, preventing the chemical branch from using
$[M/\mathrm{H}]$ as a proxy for individual elemental abundances. This architectural
choice is our primary mechanism for enforcing physically correct label--flux correlations.
We note that this state vector does not directly correspond to any named physical quantity;
it is a learned low-dimensional representation of the thermodynamic context in which
atomic line formation occurs.

\subsubsection{Chemical Expert Branch}
Each atomic species $X$ has its own dedicated expert block. The block concatenates the
thermodynamic state vector with a single scalar $[X/\mathrm{H}]$ and passes the result
through a width-adaptive expert network:
\begin{equation}
W_X = \max\!\bigl(64, \min(256, N^\mathrm{active}_X)\bigr),
\end{equation}
where $N^\mathrm{active}_X$ is the number of active pixels in that element's mask. The
expert then applies two dense layers of widths $W_X$ and $W_X/2$ before the sparse spectral
window projection.
We define the spectral windows using ASPCAP's \texttt{.filt} files
\citep{GarciaperezASPCAP2016,Abdurrouf2022},
which provide continuous-valued significance weights for each pixel--species pair rather
than binary on/off masks. The \texttt{SparseProjector} layer scatters the expert's local
optical depth prediction into the full 8575-pixel array, weighting each pixel by its
species-specific significance. Pixels with zero weight are structurally excluded from the
output:
\begin{equation}
\tau_X(\lambda_i) = w_{X,i} \, f_X\bigl(\mathbf{s}, [X/\mathrm{H}]\bigr),
\label{eq:mask}
\end{equation}
where $w_{X,i}$ is the ASPCAP window weight for species $X$ at pixel $i$,
$\mathbf{s}$ is the thermodynamic state vector, and $f_X$ is the expert sub-network.

This weighting constraint is mathematically hard: gradients cannot propagate through
zero-weight pixels, so the network is provably unable to attribute spectral signal at a
non-permitted wavelength to species $X$.

The repository currently keeps two forward architectures side by side. The full Prism
model is a two-input masked network using the per-label \texttt{bad\_mask}, while the
single-input structured ablation (\texttt{build\_final\_emulator}) uses the same sparse
atomic projectors but feeds \emph{the full input vector} into the continuum branch, the
molecular branch, the thermodynamic state branch, and the expert multilayer perceptrons.
It therefore removes the selective routing and gradient-gating logic of Prism while
retaining the final masked atomic projection. The evaluation notebook loads these as
distinct checkpoints: \texttt{APOGEE\_SGEN.keras} for Prism and
\texttt{SPECTROGRAM\_GENERATOR.keras} for Ablation Prism.

\subsection{Training}
\label{subsec:training}

We train the model end-to-end using the Adam optimizer with gradient clipping
($\|\nabla\|_{\max} = 1.0$) and an initial learning rate of $10^{-3}$. In the current
trainer this learning rate is wrapped in a cosine-decay-restarts schedule with
\texttt{first\_decay\_steps}=5000, $t_\mathrm{mul}=2.0$, $m_\mathrm{mul}=0.85$, and
$\alpha=10^{-6}$. We partition the APOGEE dataset into training (120,000 stars),
validation (20,000 stars), and held-out evaluation stars starting at index 140,000. The
active TFRecord pipeline stores spectra as 3-channel tensors
$(F,\mathit{ivar},M_\mathrm{pix})$, where $M_\mathrm{pix}$ is the per-star pixel mask
constructed from the imputation flags.

The current training notebook runs for up to 150 epochs with batch size 256, using
\texttt{ModelCheckpoint} and \texttt{EarlyStopping} (patience 20,
\texttt{min\_delta}=10^{-7}). In its present executed path, the notebook instantiates
\texttt{build\_final\_emulator}---the newer Ablation Prism baseline with masked atomic
projection but full-input routing---rather than the two-input gated Prism model. The full
Prism builder remains implemented in the same notebook and is the architecture loaded as
\texttt{APOGEE\_SGEN.keras} by the evaluation code. The other comparison models used in the
evaluation notebook---The Payne, Expanded Payne, Ablation Prism, and RegPayne---are
presently treated as fixed pre-trained checkpoints produced earlier with the same masked
spectral-pixel loss, rather than being retrained inside the active
\texttt{model-trainer.ipynb}. Relative to earlier versions of the repository, this
constitutes an optimizer upgrade: the active training path uses scheduled cosine restarts
and explicit gradient clipping instead of the older plateau-driven learning-rate reduction
workflow. At the same time, the full two-input Prism model has not undergone the same
breadth of hyperparameter fine-tuning as some baselines. This is primarily a compute
limitation, not a statement that Prism was already exhaustively optimized; the current
reported Prism results should therefore be interpreted as the performance of a
compute-constrained but fully functional architecture snapshot.

%------------------------------------------------------------
\section{Loss and Evaluation Metrics}
\label{sec:loss}
%------------------------------------------------------------

\subsection{Training Loss: Spectral Pixel Loss}
\label{subsec:spl_loss}

The active repository uses a masked spectral pixel loss. For each star we store the target
tensor as $(F_{\mathrm{obs}}, \mathit{ivar}, M_\mathrm{pix})$, where $M_\mathrm{pix}$ is a
binary or continuous keep-mask that suppresses wavelengths contaminated by imputed labels.
The loss is
\begin{equation}
\mathcal{L}_\mathrm{masked} = \sum_{i}
  M_{\mathrm{pix},i}\,w_i \bigl(F_{\mathrm{pred},i} - F_{\mathrm{obs},i}\bigr)^2,
\label{eq:spl_loss}
\end{equation}
where the per-pixel weight is
\begin{equation}
w_i = \frac{10^{-5}}{1/\mathit{ivar}_i + 10^{-5}},
\label{eq:spl_weight}
\end{equation}
and $\mathit{ivar}_i$ is the APOGEE inverse variance at pixel $i$. This weight smoothly
transitions from zero (where $\mathit{ivar} = 0$, i.e., sensor gaps) to unity (where
signal-to-noise is high), and the $10^{-5}$ floor term prevents numerical instabilities.
We additionally zero out pixels with observed flux below a bad-pixel cutoff
($F_\mathrm{obs} < 10^{-3}$) to exclude defective pixels.

\subsection{Monitoring Metric: Reduced $\chi^2$}
\label{subsec:chi2_metric}

During training, a Keras metric computes a batch-level estimate of the reduced
$\chi^2$:
\begin{equation}
\hat{\chi}^2_\mathrm{red} = \frac{1}{\nu} \sum_{i}
  M_{\mathrm{pix},i}\,\mathit{ivar}_i \, \bigl(F_{\mathrm{pred},i} - F_{\mathrm{obs},i}\bigr)^2,
\label{eq:chi2_batch}
\end{equation}
where $\nu$ is the number of valid pixels per star and we use raw
$\mathit{ivar}$ (not the scaled weight of Equation~\ref{eq:spl_weight}). This metric
provides physically interpretable goodness-of-fit during training without affecting
gradient updates.

\subsection{Test-time Evaluation: Reduced $\chi^2$}
\label{subsec:chi2_full}

The current evaluation notebook no longer reports a naive full-spectrum statistic on all
held-out stars. Instead, it uses an imputation-aware masked reduced $\chi^2$ designed to
avoid scoring models against pixels whose target labels were themselves imputed. This is
necessary because naive evaluation strongly penalizes stars with filled-in abundances:
in the current notebook, the median $\chi^2$ ratio between imputed and non-imputed stars
reaches 16.5 for C, 18.6 for O, 19.4 for Fe, and nearly 20 for Mg and Si.

We therefore define the per-star metric as
\begin{equation}
\chi^2_{\mathrm{red},\star} = \frac{1}{\nu_\star} \sum_{i=1}^{N_\mathrm{pix}}
  M_{\star i}\,
  \frac{\bigl(F_{\mathrm{pred},i} - F_{\mathrm{obs},i}\bigr)^2}
       {\sigma^2_{\mathrm{instr},i}  + 10^{-5}},
\label{eq:chi2_full}
\end{equation}
where $\sigma^2_{\mathrm{instr},i}=1/\mathit{ivar}_i$ and $M_{\star i}$ is a per-star
evaluation mask. Stars with any imputed core, CNO, or $[\mathrm{Fe/H}]$ labels are removed
entirely; for retained stars, if a trace abundance was imputed, then only that element's
sensitive windows are masked out using the APOGEE spectral masks. The degrees of freedom
$\nu_\star$ equal the number of retained pixels after masking. This protocol preserves a
fair comparison across models while preventing imputed labels from dominating the reported
forward-fit statistic.

%------------------------------------------------------------
\section{Validation}
\label{sec:validation}
%------------------------------------------------------------

We evaluate Prism on the held-out DR17 test block beginning at index 140,000. The raw
split contains 9,966 stars, but the active masked protocol drops 1,028 stars with imputed
core, CNO, or iron labels, leaving 8,938 stars for the main forward-model comparison.
We compare Prism against four pre-trained baselines that are still part of the repository:
The Payne, Expanded Payne, Ablation Prism, and RegPayne. These baselines were not removed
from the project; rather, they were already trained to completion with the masked spectral
pixel loss and are now loaded as fixed checkpoints by the evaluation notebook.

\subsection{Reduced $\chi^2$ Analysis}
\label{subsec:chi2_results}

Table~\ref{tab:chi2} reports the masked test-set statistic defined in
Equation~(\ref{eq:chi2_full}). Under this stricter protocol, all five models cluster
fairly closely in residual quality. Prism remains competitive with a median
$\chi^2_\mathrm{red}=1.718$, but it is no longer the lowest-residual model in the current
repository snapshot: The Payne and RegPayne achieve slightly lower medians of 1.683 and
1.687, respectively. The important updated interpretation is therefore a trade-off:
Prism delivers exact spectral locality and zero leakage, while several baselines achieve
slightly smaller masked residuals.

\begin{deluxetable}{lrrr}
\tablecaption{Masked reduced $\chi^2$ on the retained held-out test stars \label{tab:chi2}}
\tablewidth{0pt}
\tablehead{
  \colhead{Model} &
  \colhead{Median $\chi^2_\mathrm{red}$} &
  \colhead{Mean} &
  \colhead{95th \%}
}
\startdata
The Payne      & 1.683 & 1.890 & 3.348 \\
RegPayne       & 1.687 & 1.893 & 3.332 \\
Ablation Prism & 1.703 & 1.919 & 3.405 \\
Prism          & 1.718 & 1.941 & 3.482 \\
Expanded Payne & 1.820 & 2.114 & 3.759
\enddata

\tablecomments{All models are evaluated with the same imputation-aware masking protocol on
the retained 8,938-star subset of the held-out test block. The dropped 1,028 stars have
imputed core, CNO, or iron labels; retained stars with imputed trace abundances are scored
only on unmasked windows.}
\end{deluxetable}

We further examine $\chi^2_\mathrm{red}$ as a function of $T_\mathrm{eff}$,
$\log g$, $[M/\mathrm{H}]$, and SNR by computing binned medians across all five models
(Figures~\ref{fig:chi2_vs_teff}--\ref{fig:chi2_vs_snr}). The main qualitative trends are
shared across architectures: performance is strongest in the well-populated giant regime,
remains competitive through most of the metallicity range, and degrades sharply for hot
stars that lie well outside the dominant training distribution. Equally notable, however,
is that the binned $\chi^2$ curves in the \texttt{results/} directory have nearly the same
shape for all five models and differ mainly by modest vertical offsets. This suggests that
the models are all encountering the same underlying difficulty structure in APOGEE
spectra, and that the remaining performance gap may be governed more by the survey's
intrinsic information and noise floor than by large untapped architectural gains.

\subsection{$\chi^2$ by Stellar Type}
\label{subsec:chi2_stellar}

To probe performance across astrophysically distinct populations, we partition the test set
into six categories and compute per-group statistics (Table~\ref{tab:chi2_stellar} and
Figure~\ref{fig:chi2_distributions}).

\begin{deluxetable}{lrrrrrr}
\tablecaption{Median reduced $\chi^2_\mathrm{red}$ by stellar type \label{tab:chi2_stellar}}
\tablewidth{0pt}
\tablehead{
  \colhead{Group} &
  \colhead{$N$} &
  \colhead{Prism} &
  \colhead{Payne} &
  \colhead{Exp.P} &
  \colhead{AbP} &
  \colhead{RegP}
}
\startdata
RGB Tip          & 907  & 2.714 & 2.593 & 2.982 & 2.659 & 2.610 \\
RGB Giants       & 6485 & 1.662 & 1.627 & 1.744 & 1.646 & 1.632 \\
Subgiants        & 654  & 1.945 & 1.784 & 1.997 & 2.158 & 1.853 \\
Dwarfs           & 338  & 1.996 & 1.858 & 2.174 & 2.255 & 1.902 \\
Metal-poor Giants& 272  & 1.962 & 2.003 & 2.228 & 1.980 & 2.060 \\
Hot Stars        & 66   & 11.647 & 8.877 & 5.639 & 16.894 & 7.100 \\
\enddata
\tablecomments{Group definitions --- RGB Tip: $T_\mathrm{eff}<4000$\,K, $\log g<1.5$;
RGB Giants: $4000<T_\mathrm{eff}<5000$\,K, $1.5<\log g<3.5$;
Subgiants: $4500<T_\mathrm{eff}<6500$\,K, $3.5<\log g<4.2$;
Dwarfs: $\log g>4.2$; Metal-poor Giants: $[\mathrm{Fe/H}]<-1$, $\log g<3.5$;
Hot Stars: $T_\mathrm{eff}>6500$\,K.
Exp.P = Expanded Payne; AbP = Ablation Prism; RegP = RegPayne.
Prism no longer leads every regime. Its clearest forward-fit advantage in the current
evaluation appears for metal-poor giants, while RGB giants favor Payne-family baselines
slightly and hot stars are difficult for all models.}
\end{deluxetable}

\begin{figure}
\includegraphics[width=\textwidth]{Figures/chi2_distribution_by_type.png}
\caption{Distributions of masked reduced $\chi^2_\mathrm{red}$ by stellar type, showing
Prism (filled blue histogram) against the baseline medians (dashed vertical lines). The
main updated pattern is heterogeneous rather than uniform dominance: Prism is especially
competitive in metal-poor giants, remains close to the best model in the RGB giant regime,
and deteriorates strongly for hot stars, where all methods enter a low-information
near-infrared regime with weak metal features, dissociated molecules, and often stronger
rotational broadening.
\label{fig:chi2_distributions}}
\end{figure}

\subsection{Information Leakage Test}
\label{subsec:leakage}

An information-leakage test quantifies how much gradient power an element produces at
wavelengths where it has no known spectral lines. For each of the 14 atomic species, we
compute the sensitivity gradient $\partial F / \partial [X/\mathrm{H}]$ via forward-mode
automatic differentiation over the full test set and partition pixels into \emph{allowed}
(non-zero ASPCAP window weight) and \emph{forbidden} (zero weight) regions. The leakage
ratio is defined as $\bar{|\nabla|}_\mathrm{forbidden} \,/\,\bar{|\nabla|}_\mathrm{allowed}$;
a value well below unity indicates that the model concentrates its sensitivity on the
correct spectral windows, while a value above unity suggests the model is attributing more
spectral signal to forbidden pixels than to the element's own transitions---indicative of
abundance entanglement.

Prism achieves leakage ratios of exactly 0.000 across all 14 elements
(Table~\ref{tab:leakage}), as guaranteed by its hard masking architecture. We note that
this result is a direct consequence of the zero-weight constraint in
Equation~(\ref{eq:mask}) rather than a learned property: gradients cannot propagate through
zero-weight pixels by construction, so the zero-leakage result verifies that the
architectural constraint operates as intended rather than providing independent evidence of
physics learning.

The repository's current leakage outputs show a starker split between hard and soft
constraints than earlier drafts suggested. Ablation Prism does not consistently recover
near-local behavior: several elements exceed unity, including Fe, Si, Al, Ni, Cr, V, and
Co, with V reaching 2.44. The Payne-family baselines also retain broad non-zero leakage,
typically at the 0.1--1.2 level depending on species. The robust statement supported by
the present notebook is therefore narrower but clearer: exact locality requires more than a
masked final projector. Ablation Prism still leaks because the full label vector is fed
through shared global branches, allowing abundance information to influence flux outside the
atomic expert pathways. In the current repository, only the full two-input Prism routing
scheme---sparse atomic projection plus restricted/gated label flow---guarantees exact
locality across all 14 atomic species.

\begin{deluxetable}{lrrrrr}
\tablecaption{Information-leakage ratios ($\bar{|\nabla|}_\mathrm{forb.}/\bar{|\nabla|}_\mathrm{allow.}$) for all 14 atomic species \label{tab:leakage}}
\tablewidth{0pt}
\tablehead{
  \colhead{Element} &
  \colhead{Prism} &
  \colhead{AbP} &
  \colhead{RegP} &
  \colhead{Exp.P} &
  \colhead{Payne}
}
\startdata
Fe  & 0.000 & \textbf{1.474} & \textbf{1.051} & 0.665 & \textbf{1.061} \\
Mg  & 0.000 & 0.426 & 0.458 & 0.275 & 0.540 \\
Si  & 0.000 & \textbf{1.219} & 0.632 & 0.933 & 0.621 \\
Ca  & 0.000 & 0.377 & 0.229 & \textbf{1.187} & 0.237 \\
Ti  & 0.000 & 0.515 & 0.200 & 0.530 & 0.203 \\
S   & 0.000 & 0.366 & 0.374 & \textbf{1.083} & 0.462 \\
Al  & 0.000 & \textbf{1.019} & 0.418 & 0.884 & 0.532 \\
Mn  & 0.000 & 0.653 & 0.289 & 0.643 & 0.304 \\
Ni  & 0.000 & \textbf{1.180} & 0.602 & 0.996 & 0.569 \\
Cr  & 0.000 & \textbf{1.181} & 0.571 & 0.960 & 0.595 \\
K   & 0.000 & 0.285 & 0.122 & 0.797 & 0.134 \\
Na  & 0.000 & 0.385 & 0.181 & 0.834 & 0.197 \\
V   & 0.000 & \textbf{2.444} & 0.630 & 0.865 & 0.636 \\
Co  & 0.000 & \textbf{1.110} & 0.555 & \textbf{1.167} & 0.571 \\
\enddata
\tablecomments{AbP = Ablation Prism; RegP = RegPayne; Exp.P = Expanded Payne.
Values $> 1$ (bold) indicate more gradient power at forbidden pixels than at the
element's own transitions, a signature of abundance entanglement. Prism's zero leakage
is guaranteed by its hard masking architecture; every comparison model retains non-zero
cross-window sensitivity in the current notebook outputs.}
\end{deluxetable}

\subsection{Dependence on Stellar Parameters}
\label{subsec:param_dependence}

We examine model performance as a function of each key stellar parameter across the full
test set.

\begin{figure}
\includegraphics[width=0.49\textwidth]{Figures/median_chi2_vs_Teff_masked.png}
\includegraphics[width=0.49\textwidth]{Figures/median_chi2_vs_log_g_masked.png}
\caption{Left: Median masked reduced $\chi^2_\mathrm{red}$ versus effective temperature.
All models degrade outside the dense training manifold, especially toward hot stars and
the coolest molecularly complex giants. Right: Median masked $\chi^2_\mathrm{red}$
versus surface gravity. The main giant branch remains the easiest regime, while sparse
extremes in $\log g$ show larger model-to-model scatter. The qualitative similarity of the
curves across architectures suggests a common APOGEE-driven difficulty profile rather than
fundamentally different failure modes for each model.
\label{fig:chi2_vs_teff}\label{fig:chi2_vs_logg}}
\end{figure}

\begin{figure}
\includegraphics[width=0.49\textwidth]{Figures/median_chi2_vs_M_H_masked.png}
\includegraphics[width=0.49\textwidth]{Figures/median_chi2_vs_SNR_masked.png}
\caption{Left: Median masked $\chi^2_\mathrm{red}$ versus metallicity. Prism is most
competitive in the metal-poor giant regime, while the model ranking is otherwise fairly
tight through the bulk of the metallicity distribution. Right: Median masked
$\chi^2_\mathrm{red}$ versus SNR. Increasing SNR improves all models, but does not remove
the structural trade-off between slightly lower residuals in some baselines and exact
locality in Prism. The near-parallel profile shapes across models are consistent with all
architectures approaching the practical information limit of APOGEE-quality spectra in this
label space.
\label{fig:chi2_vs_mh}\label{fig:chi2_vs_snr}}
\end{figure}


%------------------------------------------------------------
\section{Comparison with Prior Work}
\label{sec:comparison}
%------------------------------------------------------------

\subsection{The Cannon and The Cannon 2}
\label{subsec:cannon}

The Cannon \citep{Ness2015Cannon} and The Cannon 2 \citep{Casey2016Cannon2} are
data-driven models that fit a polynomial (quadratic) generative function relating stellar
labels to spectral fluxes. The Cannon 2 uses $\ell_1$ regularization to enforce sparsity
in the model's spectral derivatives, which---like our spectral masks---prevents label
coefficients from acquiring unphysical non-zero values at all wavelengths. Our model shares
this philosophy of enforcing physical sparsity but improves upon it in two respects: first,
the sparsity in Prism is mathematically exact (hard masking) rather than approximate (soft
regularization); second, our non-linear neural architecture can capture higher-order
interactions between labels and spectra that a quadratic model cannot represent.

\subsection{The Payne}
\label{subsec:payne}

The Payne \citep{Ting2019Payne} is a multi-layer perceptron emulator trained on 2000
synthetic spectrograms with 7500 pixels, achieving a root-mean-square error of $10^{-3}$
on synthetic validation data. Because it is trained exclusively on synthetic spectra, it
is an \emph{ab initio} model: it can, in principle, return labels on any physical scale,
but it encodes all the approximations of the theoretical atmosphere grid. By contrast,
our model is trained on real APOGEE spectra and therefore learns the empirical noise floor
and calibration characteristics of the instrument.

The repository includes a Payne-architecture baseline evaluated on the same held-out DR17
split as Prism. In the current imputation-aware protocol, Payne attains a median masked
$\chi^2_\mathrm{red}=1.683$, slightly below Prism's 1.718. This means the present
comparison is not a simple win on residuals alone. Instead, Prism's distinguishing
contribution relative to Payne is structural: it achieves exact zero leakage, whereas
Payne retains non-zero forbidden-window sensitivity for every atomic species tested.

\subsection{TransformerPayne}
\label{subsec:transformerpayne}

TransformerPayne \citep{TransformerPayne2025} employs large transformer attention
mechanisms to capture long-range correlations across the full spectrum, again trained on
synthetic grids, and reports an MAE of $\sim 10^{-3}$ on a held-out synthetic dataset.
While this represents excellent performance on synthetic data, it requires substantially
more computational resources than Prism and does not include real-data validation.
Our approach achieves comparable fitting accuracy on real observations with lower
computational cost and with physically enforced interpretability.

\subsection{Baseline Comparison: Expanded Payne, Ablation Prism, and RegPayne}
\label{subsec:baselines}

In addition to the original Payne architecture, the repository keeps three further
comparison models trained on the same split. These baselines remain part of the project,
but the active trainer notebook no longer re-runs them because their masked-loss training
has already been completed and the evaluator now loads the saved checkpoints directly.
This comparison should also be read with one practical caveat: Prism was not subjected to
an equally broad hyperparameter search, largely because available compute was concentrated
on stabilizing the masked architecture itself rather than running extensive sweeps.
Accordingly, the baseline comparison is informative about the current repository state, but
it should not be over-interpreted as the final performance ceiling of the Prism design:
\begin{enumerate}
\item \textbf{Expanded Payne} --- a wider multi-layer perceptron with larger hidden layers,
      increasing parameter capacity without adding physics constraints;
\item \textbf{Ablation Prism} --- a single-input Prism variant that retains the sparse
      atomic projection masks but feeds the full label vector into the continuum,
      molecular, state, and expert networks, isolating the effect of Prism's selective
      routing and gradient-gating logic rather than removing masking entirely;
\item \textbf{RegPayne} --- a Payne-like model trained with a two-phase Jacobian penalty
      schedule (phase 1: 5 epochs, $\lambda_J = 0$; phase 2: full training,
      $\lambda_J = 400$) using the \texttt{build\_disentangled\_model} framework.
\end{enumerate}

Results are summarized in Tables~\ref{tab:chi2} and \ref{tab:chi2_stellar}. Several
observations are noteworthy. First, the masked residual ranking is now tight:
Payne (1.683), RegPayne (1.687), Ablation Prism (1.703), and Prism (1.718) all sit within
0.035 in median $\chi^2_\mathrm{red}$, while Expanded Payne trails slightly at 1.820.
Second, these small residual differences should be interpreted alongside the leakage test:
Prism is the only model with exactly zero forbidden-window response for every element.
Third, Ablation Prism no longer supports a strong claim that branch structure alone
recovers near-perfect disentanglement; its current leakage ratios exceed unity for several
species, showing that retaining masked atomic windows is not sufficient when the full label
vector is broadcast through shared branches. Fourth, the metal-poor giant population
remains the clearest niche where Prism's hard locality appears beneficial, with the best
median $\chi^2_\mathrm{red}$ in that regime among the models tested.

\subsection{Summary of Comparison}

Table~\ref{tab:comparison} summarizes the key distinctions between Prism and
representative prior work.

\begin{deluxetable}{lcccc}
\tablecaption{Comparison of stellar spectral emulators \label{tab:comparison}}
\tablewidth{0pt}
\tablehead{
  \colhead{Model} &
  \colhead{Training data} &
  \colhead{Disentangle} &
  \colhead{Interp.} &
  \colhead{Med.\ $\chi^2_\mathrm{red}$}
}
\startdata
The Cannon 2     & Real (DR12)  & Soft ($\ell_1$)  & Yes & ---   \\
The Payne        & Synth.\ grid & None             & No  & ---   \\
TransformerPayne & Synth.\ grid & None             & No  & ---   \\
\hline
Payne (ours)     & Real (DR17)  & None             & No  & \textbf{1.683} \\
Expanded Payne   & Real (DR17)  & None             & No  & 1.820 \\
Ablation Prism   & Real (DR17)  & Masked output only & Partial & 1.703 \\
RegPayne         & Real (DR17)  & Soft (Jacobian)  & Partial & 1.687 \\
\textbf{Prism}   & Real (DR17)  & Hard (mask + routing) & Yes & 1.718
\enddata
\tablecomments{Top section: published models (qualitative comparison only; dashes indicate
no reported $\chi^2_\mathrm{red}$ on real APOGEE data). Bottom section: baselines trained
on the same DR17 data split and evaluated with the masked held-out protocol. The best
current median $\chi^2_\mathrm{red}$ belongs to Payne, but Prism is uniquely
distinguished by exact zero leakage across all 14 atomic species.}
\end{deluxetable}

%------------------------------------------------------------
\section{Backward Inference on Observed Spectra}
\label{sec:backward_inference}
%------------------------------------------------------------

With Prism robustly transferring empirical patterns into its forward mapping, it naturally
supports \emph{backward inference}---recovering stellar traits from newly observed spectra
without training a separate end-to-end regressor. The maintained inverse notebook performs
fast Maximum a Posteriori (MAP) recovery by embedding the trained forward emulator inside a
two-stage optimization loop.

\subsection{Optimization Setup and Diagnostics}
\label{subsec:map_setup}
The current inverse pipeline begins with a heteroscedastic CNN warm start trained in
\texttt{CNNwarmstart.ipynb}. That network processes the observed spectrum with four 1D
convolutional blocks (32, 64, 128, and 256 channels) followed by dense layers of widths
512 and 256, and predicts both a mean label vector and a per-label variance through a
Gaussian negative log-likelihood objective. Its output supplies the initialization and
diagonal Gaussian prior for the subsequent MAP optimization.

To map an observed flux $F_\mathrm{obs}$ back to underlying labels $\theta$, we minimize a
joint objective combining the spectral reconstruction term and the warm-start prior:
\begin{equation}
\mathcal{L}_\mathrm{MAP}(\theta) = \frac{1}{2}\sum_i \mathit{ivar}_i \bigl(F_{\mathrm{obs},i} - F_{\mathrm{pred},i}(\theta)\bigr)^2 + \frac{1}{2} w_\mathrm{prior} \sum_j \left(\frac{\theta_j - \mu_j}{\sigma_j}\right)^2,
\end{equation}
where $\mu_j$ and $\sigma_j$ are the CNN-predicted means and standard deviations. The
optimization itself is FERRE-like and staged. First, a 9-dimensional core problem is fit
on the full spectrum using 1200 Adam steps. Second, the 14 trace abundances are refined in
independent one-dimensional optimizations over their masked spectral windows, using 400 Adam
steps per element. The notebook runs this procedure in batches of 32 with XLA compilation.

The current maintained notebook focuses on fast point estimation rather than calibrated
posterior uncertainty. Earlier text describing ASPCAP uncertainties as adopted confidence
bounds is therefore no longer representative of the repository; uncertainty-aware inverse
methods are being explored separately in prototype Laplace, HMC, and NUTS workflows, but
they are not the mainline result reported here.

\subsection{Validation on Stratified Test Sample}
\label{subsec:map_results}
We validated the inverse routine on a stratified sample of $N=500$ held-out stars chosen
to cover the main $T_\mathrm{eff}$, $\log g$, and $[M/\mathrm{H}]$ regimes. The full run
completed in 252.2\,s, corresponding to a median wall time of 0.31\,s per star. Table
~\ref{tab:backward_diagnostics} summarizes the resulting point-estimate diagnostics against
the APOGEE labels. Biases are near zero for most quantities, while errors are smallest for
well-constrained abundances such as Mg, Si, Ca, Mn, and Ni. The main failure mode remains
temperature recovery in the tails of the distribution, where a small number of large
outliers inflate the $T_\mathrm{eff}$ RMSE despite a much smaller robust error scale. In
particular, $T_\mathrm{eff}$ shows MAD $=10.3$\,K and IQR $=20.6$\,K but
RMSE $=218.1$\,K, while $\log g$ and $[M/\mathrm{H}]$ show the same pattern at smaller
absolute scale (MADs 0.038 and 0.013\,dex versus RMSEs 0.107 and 0.039\,dex). The large
RMSE--MAD gap implies a narrow, well-recovered core of stars plus a small set of
catastrophic failures to which the squared loss is extremely sensitive.

\begin{deluxetable}{lrrrrr}
\tablecaption{Prism Backward Inference MAP Diagnostics on $N=500$ Stratified Stars \label{tab:backward_diagnostics}}
\tablewidth{0pt}
\tablehead{
  \colhead{Label} &
  \colhead{Bias} &
  \colhead{MAD} &
  \colhead{RMSE} &
  \colhead{IQR} &
  \colhead{ASPCAP-$\sigma$}
}
\startdata
$T_\mathrm{eff}$ (K)& -1.974 & 10.284 & 218.104 & 20.629 & 8.894 \\
$\log{g}$ (dex)     & 0.002  & 0.038  & 0.107   & 0.076 & 0.025 \\
$[M/\mathrm{H}]$    & -0.002 & 0.013  & 0.039   & 0.026 & 0.007 \\
$v_\mathrm{micro}$  & -0.013 & 0.095  & 0.146   & 0.189 & 0.000 \\
$v_\mathrm{macro}$  & 0.003  & 0.020  & 0.258   & 0.041 & 0.000 \\
$v\sin i$           & -0.614 & 1.752  & 3.795   & 3.748 & 0.000 \\
$[\mathrm{C/Fe}]$   & -0.000 & 0.023  & 0.060   & 0.045 & 0.015 \\
$[\mathrm{N/Fe}]$   & 0.005  & 0.027  & 0.143   & 0.053 & 0.018 \\
$[\mathrm{O/Fe}]$   & -0.003 & 0.024  & 0.058   & 0.048 & 0.018 \\
$[\mathrm{Fe/H}]$   & -0.003 & 0.013  & 0.041   & 0.027 & 0.008 \\
$[\mathrm{Mg/Fe}]$  & -0.001 & 0.009  & 0.058   & 0.018 & 0.013 \\
$[\mathrm{Si/Fe}]$  & -0.001 & 0.012  & 0.070   & 0.023 & 0.015 \\
$[\mathrm{Ca/Fe}]$  & -0.001 & 0.013  & 0.054   & 0.027 & 0.015 \\
$[\mathrm{Ti/Fe}]$  & -0.000 & 0.021  & 0.095   & 0.042 & 0.024 \\
$[\mathrm{S/Fe}]$   & 0.002  & 0.020  & 0.078   & 0.040 & 0.044 \\
$[\mathrm{Al/Fe}]$  & -0.002 & 0.022  & 0.096   & 0.044 & 0.020 \\
$[\mathrm{Mn/Fe}]$  & 0.002  & 0.013  & 0.054   & 0.025 & 0.017 \\
$[\mathrm{Ni/Fe}]$  & -0.001 & 0.012  & 0.046   & 0.024 & 0.013 \\
\enddata
\tablecomments{Values represent point-estimate diagnostics against ASPCAP labels for the
current fast-MAP pipeline. The maintained notebook reports point estimates only; the
ASPCAP column is included as an external scale reference, not as a claimed posterior
uncertainty for Prism. Here MAD denotes the median absolute deviation from the ASPCAP
reference label.}
\end{deluxetable}

The current maintained repository does not include a synchronized cross-model inverse
benchmark analogous to the forward comparison above. Earlier placeholder tables comparing
MAP errors across baselines are therefore removed here. Cross-model inverse benchmarking is
a natural next step, but the present codebase validates only the Prism MAP pipeline
end-to-end.

%------------------------------------------------------------
\section{Discussion}
\label{sec:discussion}
%------------------------------------------------------------

We have presented a neural emulator for stellar spectra that combines the advantages of
data-driven modeling---training on real observations, robustness to noise and calibration
artifacts---with hard physical constraints that guarantee interpretable, disentangled label
sensitivities. The architectural separation of the spectrum into continuum, molecular, and
atomic contributions mirrors the physical origin of these spectral components and
provides a natural pathway for model interpretation.

The current repository points to a more nuanced trade-off than earlier drafts emphasized.
On masked forward residuals, Prism is competitive but not universally best; several
baseline checkpoints achieve slightly lower median $\chi^2_\mathrm{red}$. Its strongest
empirical niche is the metal-poor giant regime, where weak-line abundances benefit most
from strict spectral locality. The defining result is instead the leakage test:
hard masking is the only mechanism in the present codebase that delivers exact
zero forbidden-window response for all 14 atomic species.

\subsubsection{Limitations}

Several limitations of the current model merit attention. First, the thermodynamic state
branch introduces a degree of flexibility that, if unconstrained, could in principle absorb
variance attributable to chemical abundances. In practice, we manage this through the
architectural constraint that the chemical expert branch receives only the state vector
(not raw labels), but more rigorous sensitivity auditing would strengthen this claim.

Second, although our imputation procedure rigorously preserves data isolation across partitions by pre-fitting to the training set alone, it still introduces a 
correlation structure into the labels that may not completely reflect true astrophysical
covariance. This could bias the thermodynamic state vector for stars with many missing
labels. A more principled treatment, such as explicit latent variable models for missing
data, remains a natural extension.

Third, although the repository now includes a fast backward inference pipeline, it is still
primarily a point-estimation system. Calibrated uncertainty quantification remains under
active development in separate Laplace, HMC, and NUTS prototypes and is not yet part of
the mainline workflow described in this paper.

Fourth, being a deterministic model, Prism does not reproduce the stochastic stellar
variance that would exist among stars sharing the same parameters, such as variance arising
through non-LTE physics or magnetic activity. This limits its utility as a synthetic data
generator for population-level applications.

\subsubsection{Performance in Temperature Extremes}

Systematic degradation in cool regions ($T_\mathrm{eff} < 4500$\,K) is the most notable
limitation. This stems from two factors: (i) the training sample contains relatively few
extremely cool stars (${\sim}$30\% of the dataset falls below 4500\,K), so the model has
limited exposure to this regime, and (ii) molecular physics and line lists are incomplete,
meaning even perfect physical constraints cannot fully compensate for missing spectroscopic
data. This is not a deficiency of our approach per se---all baselines show similar
degradation---but rather a fundamental limitation of the underlying data. The current
evaluation notebooks also show that hot stars ($T_\mathrm{eff}>6500$\,K) are a major
failure mode for all models, including Prism. This is not only an out-of-distribution
issue. In the APOGEE $H$ band, hotter atmospheres ionize many metals and dissociate the
molecules that dominate cool-star diagnostics, so the spectrum intrinsically loses a large
fraction of its abundance-sensitive structure. The remaining lines are often shallower and,
for many hot stars, additionally smeared by larger projected rotation. From a physical
standpoint, the inverse and forward problems are therefore both less identifiable in this
regime: there is simply less chemically specific information available in the APOGEE
wavelength range.

More broadly, the near-identical shapes of the binned $\chi^2$ profiles across Prism and
all baselines suggest that the dominant residual structure is shared. This points to a
practical ceiling set by APOGEE's wavelength coverage, label noise, normalization systematics,
and intrinsic information content, with architecture primarily affecting locality and
calibration rather than radically changing the global error landscape.

%------------------------------------------------------------
\section{Future Work}
\label{sec:future}
%------------------------------------------------------------

We outline the following planned extensions and validation experiments.

\paragraph{External validation via classifier.}
We will train a discriminative neural network to distinguish real APOGEE spectra from
Prism-generated ones. If the classifier's accuracy approaches chance level (50\%), it will
confirm that generated and real spectra are statistically indistinguishable.

\paragraph{Synthetic dataset generator.}
We will release a public Python class that allows users to specify a parameter space and
generate arbitrary numbers of synthetic spectra for downstream applications such as
training data-driven classifiers or testing spectral analysis pipelines.

\paragraph{Inverse uncertainty quantification.}
The forward model now supports fast MAP inversion, but the next major step is a maintained
uncertainty-aware inverse stack. Current prototypes include Laplace approximations and
sampling-based HMC/NUTS approaches built around the same forward emulator.

\paragraph{Extension to other surveys.}
While trained on APOGEE H-band, the architecture is survey-agnostic. We envision
extensions to visible-wavelength surveys (GALAH, LAMOST) and to other infrared bands,
provided appropriate training data and spectral masks are available.

%------------------------------------------------------------
\section{Conclusion}
\label{sec:conclusion}
%------------------------------------------------------------

We have constructed a physics-constrained generative neural emulator for APOGEE DR17
stellar spectra. We train Prism on 149,966 real observations and, under the current
imputation-aware masked evaluation, obtain a median reduced $\chi^2_\mathrm{red}=1.718$
on the retained 8,938-star held-out subset. This is competitive rather than dominant in
residual terms, but it is paired with a stronger structural guarantee than any baseline in
the repository: exact zero leakage outside each element's permitted spectral windows.

The present codebase therefore supports a sharpened conclusion. Hard spectral locality does
not automatically minimize masked $\chi^2$ in every regime, but it does enforce complete
disentanglement and appears especially valuable in weak-line settings such as metal-poor
giants. The same repository also now demonstrates a practical inverse workflow: a
heteroscedastic CNN warm start plus staged XLA-accelerated MAP optimization recovers labels
for a 500-star stratified sample in 0.31\,s per star on median.

Prism represents a step toward spectral emulators that are simultaneously competitive,
physically transparent, and faithful to the observational properties of the instruments
that collect the data.

All code and trained model weights are available in our project repository, with
comprehensive documentation and usage examples. We welcome community contributions and
discussion via GitHub.

%------------------------------------------------------------

\section{Author Contributions}
Aneesh Shastri conceived the model architecture, collected and preprocessed the data,
performed all training and validation experiments, produced all figures, and wrote the
manuscript.


\begin{acknowledgments}
%------------------------------------------------------------

Thanks to Prof. Mayukh Pahari (Physics Department, IITH) for guidance throughout this
project.

Thanks to Anurag Baruah for assistance in formalizing results and providing essential
guidance on research paper writing.

Thanks to friends and family for support throughout the paper-drafting process. We
gratefully acknowledge Cepheid, IITH for providing the opportunity to go through with
this project.

This project made use of data from the SDSS-IV APOGEE survey. Funding for
the Sloan Digital Sky Survey IV has been provided by the Alfred P.\ Sloan Foundation,
the U.S.\ Department of Energy Office of Science, and the Participating Institutions.
SDSS-IV acknowledges support and resources from the Center for High-Performance Computing
at the University of Utah. The SDSS web site is \texttt{www.sdss.org}.

This research made use of Astropy, a community-developed core Python package for
Astronomy \citep{Astropy2013}; NumPy \citep{vanderWalt2011}; SciPy \citep{Jones2001};
Matplotlib \citep{Hunter2007}; and scikit-learn \citep{Pedregosa2011}.
\end{acknowledgments}
%------------------------------------------------------------
\begin{thebibliography}{}

\bibitem[Abdurrouf et al.(2022)]{Abdurrouf2022}
Abdurrouf, et al.\ 2022, ApJS, 259, 35

\bibitem[Astropy Collaboration et al.(2013)]{Astropy2013}
Astropy Collaboration, Robitaille, T.\ P., Tollerud, E.\ J., et al.\ 2013, A\&A, 558, AA33

\bibitem[Bland-Hawthorn \& Gerhard(2016)]{BlandHawthorn2016}
Bland-Hawthorn, J., \& Gerhard, O.\ 2016, ARA\&A, 54, 529

\bibitem[Casey et al.(2016)]{Casey2016Cannon2}
Casey, A.\ R., Hogg, D.\ W., Ness, M., Rix, H.-W., Ho, A.\ Y.\ Q., \& Gilmore, G.\
2016, arXiv:1603.03040

\bibitem[Freeman \& Bland-Hawthorn(2002)]{Freeman2002}
Freeman, K., \& Bland-Hawthorn, J.\ 2002, ARA\&A, 40, 487

\bibitem[Garc\'ia P\'erez et al.(2016)]{GarciaperezASPCAP2016}
Garc\'ia P\'erez, A.\ E., Allende Prieto, C., Holtzman, J.\ A., et al.\ 2016,
AJ, 151, 144

\bibitem[Hunter(2007)]{Hunter2007}
Hunter, J.\ D.\ 2007, Comput.\ Sci.\ Eng., 9, 90

\bibitem[Jones et al.(2001)]{Jones2001}
Jones, E., Oliphant, E., Peterson, P., et al.\ 2001, SciPy: Open Source Scientific
Tools for Python, \url{http://www.scipy.org/}

\bibitem[Majewski et al.(2017)]{Majewski2017}
Majewski, S.\ R., Schiavon, R.\ P., Frinchaboy, P.\ M., et al.\ 2017, AJ, 154, 94

\bibitem[Ness et al.(2015)]{Ness2015Cannon}
Ness, M., Hogg, D.\ W., Rix, H.-W., Ho, A.\ Y.\ Q., \& Zasowski, G.\ 2015,
ApJ, 808, 16

\bibitem[Pedregosa et al.(2011)]{Pedregosa2011}
Pedregosa, F., Varoquaux, G., Gramfort, A., et al.\ 2011, JMLR, 12, 2825

\bibitem[Shetrone et al.(2015)]{Shetrone2015}
Shetrone, M., Bizyaev, D., Lawler, J.\ E., et al.\ 2015, ApJS, 221, 24

\bibitem[Smith et al.(2021)]{Smith2021DR17}
Smith, V.\ V., Cunha, K., Shetrone, M.\ D., et al.\ 2021, AJ, 161, 254

\bibitem[Ting et al.(2019)]{Ting2019Payne}
Ting, Y.-S., Conroy, C., Rix, H.-W., \& Cargile, P.\ 2019, ApJ, 879, 69

\bibitem[TransformerPayne(2025)]{TransformerPayne2025}
TransformerPayne Collaboration\ 2025, arXiv preprint [astro-ph.SR]

\bibitem[van der Walt et al.(2011)]{vanderWalt2011}
van der Walt, S., Colbert, S.\ C., \& Varoquaux, G.\ 2011, Comput.\ Sci.\ Eng., 13, 22

\end{thebibliography}

\end{document}
